\documentclass[10pt,a4paper]{amsart}
\usepackage[margin=19mm]{geometry}
\usepackage{amsmath,amssymb,mathtools}
\usepackage[scr=rsfs,cal=euler]{mathalfa}
\usepackage{xfrac}
\usepackage[unicode,hidelinks,bookmarks=false]{hyperref}
\newcommand{\ie}{i.e.\ }
\newcommand{\cf}{cf.\ }
\newcommand{\resp}{respectively\ }
\newcommand{\Subsection}{\subsection}
\providecommand{\nobreakdash}{\nobreak\hbox{-}}
\setlength{\emergencystretch}{2em}
\multlinegap0pt
\usepackage{bm}
\usepackage{bbm}
\usepackage{dsfont}
\usepackage{xspace}
\usepackage{cancel}
\usepackage{mathtools}
\usepackage{amsthm}
\makeatletter
\@ifundefined{thm}{\newtheorem{thm}{Thm}}{}
\@ifundefined{theorem}{\newtheorem{theorem}[thm]{Theorem}}{}
\makeatother
\title[Water Cells in Compositions of 1s and 2s]{Water Cells in Compositions of 1s and 2s}
\author{Brian Hopkins, Aram Tangboonduangjit}
\date{Source: arXiv:2412.11528v2 (CC BY 4.0)}
\begin{document}
\maketitle

\noindent\emph{Source and licence.} Water Cells in Compositions of 1s and 2s. Version arXiv:2412.11528v2 (CC BY 4.0), distributed under the Creative Commons Attribution 4.0 International licence, as recorded in the arXiv licence metadata for this exact version and shown on the abstract page https://arxiv.org/abs/2412.11528v2. Full rights notice: licenses/arxiv-2412.11528-CC-BY-4.0.txt.\par\smallskip

\noindent\textit{Editorial scope.} Counted unit: the Theorem whose source label is \texttt{wum} in the article (source0.tex lines 261--275, source label \texttt{wum}), delivered together with the article's own complete original proof of it (source0.tex lines 277--306, one \texttt{proof} environment of 30 lines). No mathematics is written, repaired or re-derived: statement and proof are the article's own text line for line. Required context retained uncounted: wcgf (lines 207--209). Every definition, display and label of those lines is retained unchanged. Omitted from this package and not redistributed: 1--206, 210--260, 276--276, 307--633. Nothing inside a retained excerpt is omitted or altered: every retained body is delivered exactly as the article's lines are written, so no omission receipt of any kind accompanies this package. Citations: the retained excerpts carry no citation command at all, so no bibliography entry of the article is transcribed and no reference is redistributed. Reference adaptation: none was needed and none was made; every reference inside the delivered unit resolves inside it, so no unresolved reference or citation is produced. Printed slips kept literal: no printed word, symbol, hypothesis or number of the counted statement or of its proof was altered. Layout and class: the article's own class is \texttt{interact} with packages including latexsym, anyfontsize, tikz, color, hyperref, url, breakurl; the wrapper is the lane's pinned \texttt{amsart} editorial wrapper at 10pt on A4, which restates the theorem-like environments the retained text needs and adds the article's own macro definitions that the retained text needs (none), with extra wrapper packages (bm, bbm, dsfont, xspace, cancel, mathtools). The delivered numbering is the unit's own. Assets: no retained span contains a figure environment, a graphics-inclusion command, a PDF-page inclusion, a table or a TikZ picture; no image file is copied into this package, the package declares no asset dependency and no image file is redistributed with it. Licence: version arXiv:2412.11528v2 of this article is distributed under the Creative Commons Attribution 4.0 International licence, as recorded in the arXiv licence metadata for this exact version and shown on the abstract page https://arxiv.org/abs/2412.11528v2. A full rights notice is delivered at licenses/arxiv-2412.11528-CC-BY-4.0.txt. Characters: every retained excerpt is pure ASCII, so the delivered engine, XeLaTeX with its default Latin Modern text font, reports no missing character.
\begin{equation}
\sum w(n,k) q^n z^k = \frac{1}{1-q} + \frac{q^2(1-zq)}{(1-q)^2(1-zq-q^2)} \label{wcgf}
\end{equation}

\begin{theorem} \label{wum}
The generating function for $w(n,0)$ is
\begin{equation}
 \sum_{n \ge 0} w(n,0) q^n = \frac{1}{1-q} + \frac{q^2}{(1-q)^2(1-q^2)} = \frac{1-q+q^3}{1 - 2 q + 2 q^3 - q^4}. \label{w0gf}
 \end{equation}

Also, for each $n \ge 0$, the following values are equal.
\begin{enumerate}
\renewcommand{\theenumi}{\alph{enumi}}
\item $w(n,0)$,
\item $\lfloor n^2/4 \rfloor + 1$,
\item $w(n-2,0) + n - 1$ (for $n \ge 2$),
\item $2w(n-1,0) - 2w(n-3,0) + w(n-4,0)$ (for $n \ge 4$).
\end{enumerate}
\end{theorem}

\begin{proof}
For the generating function, rather than work from \eqref{wcgf}, we give a combinatorial argument.  A composition in $C_{12}(n)$ with no water cells is weakly unimodal, so it has the form $(1^a, 2^b, 1^c)$ for nonnegative $a,b,c$.  One possibility is the composition $(1^n) \in C_{12}(n)$ accounted for in the generating function by the summand $1/(1-q)$.  The second summand, $q^2/((1-q)^2(1-q^2))$, counts partitions of $n$ with parts 1 and 2 where there is at least one part 2 and there are two colors of parts 1.  These correspond to compositions in $C_{12}(n)$ with no water cells and at least one part 2: any partition parts $1_1$ correspond to the initial run $1^a$ in the composition, additional parts 2 from the partition contribute to $2^b$, and any partition parts $1_2$ correspond to the terminal run $1^c$.  (For example, the partition $(2,2,1_1,1_1,1_2)$ corresponds to the composition $(1,1,2,2,1)$.)  The second generating function expression follows from algebraic manipulation.

For the direct formula (b), consider the number of parts 2, which must be adjacent.  Suppose first that $n = 2m$.  There is one composition counted by $w(n,0)$ that has no parts 2, the composition $(1^n)$.  There are $2m-1$ compositions with a single part 2, as it can be before or after any of the $2m-2$ parts 1.  There are $n-3$ compositions with the block $(2,2)$ before or after any of the $n-4$ parts 1.  This continues to the single composition $(2^m)$, giving
\[w(2m,0) = 1 + (2m-1) + (2m-3) + \cdots + 1 = m^2 + 1.\]
The analogous sum for the case $n = 2m+1$ is
\[w(2m+1,0) = 1 + 2m + (2m-2) + \cdots + 2 = m^2 + m + 1\]
and, in both cases, these match $\lfloor n^2/4 \rfloor + 1$.

For the nonhomogeneous recurrence (c), we describe how to construct $W(n,0)$ from most of $W(n-2,0)$ and some other compositions.  We actually show 
\[w(n,0) = (w(n-2,0) - 1) + (n-1) + 1.\]
For each of the $w(n-2,0) - 1$ compositions of $W(n-2,0)$ that contain at least one part 2 in a single run, add another part 2 to that run.  There are $n-1$ compositions in $W(n,0)$ with a single part 2 in every possible positions before and after the $n-2$ parts 1.  Finally, $(1^n) \in W(n,0)$.  In the reverse direction, from the compositions in $W(n,0)$ with two or more parts 2, remove one of them to make a composition in $W(n-2,0)$; this leaves $n-1$ compositions in $W(n,0)$ with a single part 2 and also $(1^n)$.

The linear recurrence (d) follows from the generating function \eqref{w0gf}, but we give a combinatorial proof, establishing the bijection
\[W(n,0) \cup 2W(n-3,0) \cong 2W(n-1,0) \cup W(n-4,0)\]
where the coefficient 2 means two copies of a set.  The explanation of the bijection is simplified by defining the increasable subset $W^i(n,0)$ of $W(n,0)$ as $(1^n)$ and the compositions whose last two parts are $(2,1)$, i.e., the compositions where the last part can be increased by one without creating a water cell.  Let $W^j(n,0) = W(n,0) \setminus W^i(n,0)$.  We show 
\begin{gather*}
W(n,0) \cong W(n-1,0) \cup W^i(n-1,0), \\
2W(n-3,0) \cong\ W^j(n-1,0) \cup W(n-4,0)
\end{gather*}
from which the claim follows.

For the first bijection, decrease the last part of each composition in $W(n,0)$ by one.  Those with last part 1 are mapped bijectively to $W(n-1,0)$.  The compositions in $W(n,0)$ with last part 2 are mapped into another $W(n-1,0)$ set, namely the compositions ending in $(2,1)$ and $(1^{n-1})$, i.e., $W^i(n-1,0)$.

For the reverse map of the first bijection, add a part 1 at the end of the compositions in $W(n-1,0)$ to obtain the compositions of $W(n,0)$ with last part 1.  Increase the last part of the compositions in $W^i(n-1,0)$ by one (allowed by definition), giving the compositions of $W(n,0)$ with last part 2.

For the second bijection, in the first set $W(n-3,0)$, for the compositions that end in 1, removing that last part establishes a bijection to $W(n-4,0)$.  For the compositions in the first set $W(n-3,0)$ that end in 2, add another part 2.  In the second set $W(n-3,0)$, send $(1^{n-3})$ to $(1^{n-3},2)$ and add the parts $(1,1)$ at the end of the other compositions.  Together these comprise the set $W^j(n-1,0)$: those ending with a run of two or more parts 2, the composition ending in a single 2, and those with at least one part 2 ending in $(1,1)$, exactly the compositions where increasing the last part yields a composition outside $C_{12}(n-1)$ or a composition in $C_{12}(n-1)$ with a water cell.

For the reverse map of the second bijection, for the compositions of $W^j(n-1,0)$ that end in $(1,1)$, removing those last two parts establishes a bijection to $W(n-3,0)$ except for $(1^{n-3})$ (since $(1^{n-1}) \in W^i(n-1,0)$).  Since no composition of $W^j(n-1,0)$ ends in $(2,1)$ (that would be increasable), the remaining compositions have last part 2, which we remove.  One of those is $(1^{n-3},2) \in W^j(n-1,0)$ whose image $(1^{n-3})$ completes the first $W(n-3,0)$.  The other compositions in $W^j(n-1,0)$ that end in 2 have additional parts 2: These must be in a run of parts 2 at the end, else there would be a water cell.  Thus the images in the second $W(n-3,0)$ comprise all the compositions in that $W(n-3,0)$ that end in 2.  Finally, add a part 1 at the end of each composition in $W(n-4,0)$ to complete the second $W(n-3,0)$ with the compositions that end in 1.
\end{proof}

\end{document}
